% Part: normal-modal-logic
% Chapter: tableaux
% Section: introduction

\documentclass[../../../include/open-logic-section]{subfiles}

\begin{document}

\olfileid{nml}{tab}{int}

\olsection{పరిచయం}

\tetoken{టాబ్లోలు}{!!^{tableau}s} అంటే కిందికి శాఖలుగా విస్తరించే
\tetoken{చిహ్నిత సూత్రాల}{!!{signed
  formula}s} వృక్షాలు. చిహ్నిత సూత్రం అనేది ఒక సత్యమూల్య
సంకేతం ($\True$ లేదా $\False$), \tetoken{ఒక వాక్యంతో}{!!a{sentence}}
కూడిన జంట:
\[
\sFmla{\True}{!A} \text{ లేదా } \sFmla{\False}{!A}.
\]
\tetoken{ఒక టాబ్లో}{!!^a{tableau}} కొన్ని \emph{పరికల్పనలతో}
మొదలవుతుంది. ఆ తరువాతి ప్రతి \tetoken{చిహ్నిత సూత్రం}{!!{signed formula}}
నిగమన నియమాల్లో ఒకదాన్ని వర్తింపజేయడం వల్ల వస్తుంది. కొన్ని
నియమాలు వృక్షపు శాఖ చివరికి ఒకటి లేదా అంతకంటే ఎక్కువ
\tetoken{చిహ్నిత సూత్రాలను}{!!{signed formula}s} చేర్చుతాయి;
మరికొన్ని రెండు కొత్త చివరలను చేర్చి రెండు శాఖలను ఏర్పరుస్తాయి.
నియమాన్ని వర్తింపజేసిన \tetoken{చిహ్నిత సూత్రంలోని}{!!{signed formula}}
సూత్రం కంటే తక్కువ సంక్లిష్టతగల
\tetoken{సూత్రం}{!!{formula}} ఉన్న \tetoken{చిహ్నిత సూత్రాలను}{!!{signed formula}s}
ఆ నియమాలు ఇస్తాయి. ఒక శాఖలో $\sFmla{\True}{!A}$,
$\sFmla{\False}{!A}$ రెండూ ఉంటే ఆ శాఖ \emph{సంవృతమైనది} అంటాం.
\tetoken{ఒక టాబ్లోలోని}{!!a{tableau}} ప్రతి శాఖ సంవృతమైతే
ఆ \tetoken{టాబ్లో}{!!{tableau}} మొత్తం సంవృతమైనది. సంవృత
\tetoken{టాబ్లో}{!!{tableau}}, \tetoken{టాబ్లోను}{!!{tableau}}
ప్రారంభించడానికి ఉపయోగించిన
\tetoken{చిహ్నిత సూత్రాల}{!!{signed formula}s} సమితి
సంతృప్తిపరచలేనిదని చూపే \tetoken{వ్యుత్పత్తి}{!!a{derivation}}.
దీని ద్వారా $\Proves$ సంబంధాన్ని ఇలా నిర్వచించవచ్చు:
$\Gamma \Proves !A$ అయితే, అప్పుడు మాత్రమే దానిలోని ఏదో
పరిమిత సమితి~$\Gamma_0 = \{!B_1, \dots, !B_n\} \subseteq \Gamma$
కోసం కింది పరికల్పనలకు సంవృత \tetoken{టాబ్లో}{!!{tableau}} ఉంటుంది:
\[
\{\sFmla{\False}{!A}, \sFmla{\True}{!B_1}, \dots, \sFmla{\True}{!B_n}\}.
\]

మోడల్ తర్కాల కోసం \tetoken{చిహ్నిత సూత్రం}{!!{signed
formula}} అనే భావనను విస్తరించడంతోపాటు,
\iftag{prvBox}{$\Box$\iftag{prvDiamond}{ మరియు
    $\Diamond$}}{$\Diamond$}ను కవర్ చేసే నియమాలనూ చేర్చాలి.
సంకేతం ($\True$ లేదా $\False$)తోపాటు, మోడల్
\tetoken{టాబ్లోలలోని}{!!{tableau}s} \tetoken{సూత్రాలకు}{!!{formula}s}
\emph{పూర్వసూచికలు}~$\sigma$ కూడా ఉంటాయి. పూర్వసూచికలు ధన
పూర్ణసంఖ్యల శూన్యం కాని అనుక్రమాలు; అంటే
$\sigma \in (\PosInt)^* \setminus \{\emptyseq\}$.
అటువంటి పూర్వసూచికలను చుట్టూ $\tuple{\ }$ లేకుండా రాస్తాం;
వాటిలోని \tetoken{మూలకాలను}{!!{element}s} వేరుచేయడానికి
$.$లను ఉపయోగిస్తాం, $,$లను కాదు. $\sigma$ ఒక పూర్వసూచిక అయితే
$\sigma.n$ అంటే $\sigma \concat \tuple{n}$; ఉదాహరణకు,
$\sigma = 1.2.1$ అయితే $\sigma.3$ అనేది $1.2.1.3$.
కాబట్టి, ఉదాహరణకు,
\[
\sFmla{\True}{\Box !A \lif !A}[1.2]
\]
అనేది \emph{పూర్వసూచిక గల \tetoken{చిహ్నిత సూత్రం}{!!{signed formula}}}
(సంక్షిప్తంగా \emph{పూర్వసూచిక గల \tetoken{సూత్రం}{!!{formula}}}).

సహజ అవగాహన ప్రకారం, పూర్వసూచిక ఒక నమూనాలోని లోకానికి
పేరు. ఆ లోకం \tetoken{ఒక టాబ్లోలోని}{!!a{tableau}} శాఖపై ఉన్న
\tetoken{సూత్రాలను}{!!{formula}s} సంతృప్తిపరచవచ్చు.
$\sigma$ ఒక లోకానికి పేరైతే, $\sigma.n$ అనేది~$\sigma$ పేరుగల
లోకం నుంచి ప్రాప్యమైన లోకానికి పేరు.

\end{document}
